%!tex
\input macros.tex


\section{Revis\~ao}


<<<<<<< HEAD
Qual \'e a ideia fundamental do projeto? Basicamente, usar ampops NE5532 em paralelo para atingir uma corrente aceit\'avel, com o fito de suportar uma carga de $8\,\Omega$ de imped\^ancia. Sabe-se que um \'unico ampop que comp\~oe o amplicador dual NE5532 \'e capaz de alimentar uma carga de $500\,\Omega$ com a tens\~ao de alimenta\c c\~ao do ampop, digamos $18.3\,\hbox{V}$, i.e., em teoria temos em \'unico ampop uma corrente de sa\'\i da de 
=======
Qual \'e a ideia fundamental do projeto? Basicamente, usar ampops NE5532 em paralelo para atingir uma corrente aceit\'avel, com o fito de suportar uma carga de $8\,\Omega$ de imped\^ancia. Sabe-se que um \'unico ampop que comp\~oe o amplicador dual NE5532 \'e capaz de alimentar uma carga de $500\,\Omega$ com a tens\~ao de alimenta\c c\~ao do ampop, digamos com uma margem de erro de $3\,\hbox{V}$, ou seja, $15.3=18.3-3\,\hbox{V}$, i.e., em teoria temos em \'unico ampop uma corrente de sa\'\i da de 
>>>>>>> 7e586d1 (run: 15.03.26 às 11:53)
$$
  {15.3\,\hbox{V}\over 500}={15.3\,\hbox{V}\over 500\sqrt2\,\Omega}\approx22\,\hbox{mA}_{\rm rms}.
$$
Se colocarmos em paralelo 64 destes ampops, temos pela lei de Kirchhoff dos n\'os, que num n\'o de conflu\^encia a corrente total \'e
$$
  I_T=64\cdot22\,\hbox{mA}_{\rm rms}\approx1.38\,\hbox{A}_{\rm rms}.
$$
Em conformidade, isso d\'a uma pot\^encia te\'orica de 
$$
  P_{\rm rms}={64\cdot15.3^2\over1000}\approx15\,\hbox{W}_{\rm rms}
$$
Conforme mencionado no artigo da Elektor.





\bye
